\begin{frame}{まとめ}{}
  \begin{cards}[2]
    \palecard{北海道の高層風の姿を捉えた}{秋にジェット最強。対流圏の風向は年中読め、破裂高度帯は春秋に読めない。}
    \palecard{意味のある統計期間は 1〜2 か月}{最細は 15 日。上空の対流圏は暦で当たらない（落下位置で 17 \%）。}
    \palecard{VOISING からは通年で海上}{10 月下旬は東北東へ 185 km、岸から 55 km。1 日ずれで 56 km 動く。}
    \palecard[wcOrange]{気候値は計画の道具}{風の固定の近似は 3--5 km で、日程のずれの 1 桁下。当日の判断は予報で。}
  \end{cards}
\end{frame}
